\documentclass[10pt,a4paper]{amsart}
\usepackage{amsmath,amssymb,amsthm,graphicx}
% Retained preamble for the Stacks Project chapter snapshot in this directory.
% The chapter's own preamble.tex is a copy of the shared Stacks Project preamble; the block
% below is sliced MECHANICALLY, line for line, from it. Nothing was added, removed or
% redefined; the wrapper class, the xr-hyper/hyperref wiring and the externaldocument list
% are omitted because the assembler supplies the document class.
%
% lines taken from preamble.tex: 3, 5, 9, 11, 12, 156, 157, 158, ... (51 lines total)
%
% package coverage restored from the chapter's own preamble (2026-10-13): xy, xy-2cell

\usepackage[all]{xy}
\xyoption{2cell}
\UseAllTwocells

\newenvironment{reference}{\comment}{\endcomment}
\newenvironment{slogan}{\comment}{\endcomment}
\newenvironment{history}{\comment}{\endcomment}
\theoremstyle{plain}
\newtheorem{theorem}[subsection]{Theorem}
\newtheorem{proposition}[subsection]{Proposition}
\newtheorem{lemma}[subsection]{Lemma}
\theoremstyle{definition}
\newtheorem{definition}[subsection]{Definition}
\newtheorem{example}[subsection]{Example}
\newtheorem{exercise}[subsection]{Exercise}
\newtheorem{situation}[subsection]{Situation}
\theoremstyle{remark}
\newtheorem{remark}[subsection]{Remark}
\newtheorem{remarks}[subsection]{Remarks}
\numberwithin{equation}{subsection}
\def\lim{\mathop{\mathrm{lim}}\nolimits}
\def\colim{\mathop{\mathrm{colim}}\nolimits}
\def\Spec{\mathop{\mathrm{Spec}}}
\def\Hom{\mathop{\mathrm{Hom}}\nolimits}
\def\Ext{\mathop{\mathrm{Ext}}\nolimits}
\def\SheafHom{\mathop{\mathcal{H}\!\mathit{om}}\nolimits}
\def\SheafExt{\mathop{\mathcal{E}\!\mathit{xt}}\nolimits}
\def\Sch{\mathit{Sch}}
\def\Mor{\mathop{\mathrm{Mor}}\nolimits}
\def\Ob{\mathop{\mathrm{Ob}}\nolimits}
\def\Sh{\mathop{\mathit{Sh}}\nolimits}
\def\NL{\mathop{N\!L}\nolimits}
\def\CH{\mathop{\mathrm{CH}}\nolimits}
\def\proetale{{pro\text{-}\acute{e}tale}}
\def\etale{{\acute{e}tale}}
\def\QCoh{\mathit{QCoh}}
\def\Ker{\mathop{\mathrm{Ker}}}
\def\Im{\mathop{\mathrm{Im}}}
\def\Coker{\mathop{\mathrm{Coker}}}
\def\Coim{\mathop{\mathrm{Coim}}}
\DeclareMathSymbol{\boxtimes}{\mathbin}{AMSa}{"02}
\def\QCohstack{\mathcal{QC}\!\mathit{oh}}
\def\Cohstack{\mathcal{C}\!\mathit{oh}}
\def\Spacesstack{\mathcal{S}\!\mathit{paces}}
\def\Quotfunctor{\mathrm{Quot}}
\def\Hilbfunctor{\mathrm{Hilb}}
\def\Curvesstack{\mathcal{C}\!\mathit{urves}}
\def\Polarizedstack{\mathcal{P}\!\mathit{olarized}}
\def\Complexesstack{\mathcal{C}\!\mathit{omplexes}}
\def\Pic{\mathop{\mathrm{Pic}}\nolimits}
\def\Picardstack{\mathcal{P}\!\mathit{ic}}
\def\Picardfunctor{\mathrm{Pic}}
\def\Deformationcategory{\mathcal{D}\!\mathit{ef}}
\usepackage{cleveref}
\crefname{originallemma}{original lemma}{original lemmas}
\Crefname{originallemma}{Original lemma}{Original lemmas}
\makeatletter
\newlabel{lemma-dg-hull}{{0FSK}{source}{}{sourcepointer-1}{}}
\newlabel{lemma-dg-hull@cref}{{[originallemma][0][]0FSK}{[1][0][]source}{}{}{}}
\makeatother
\crefname{originallemma}{original lemma}{original lemmas}
\Crefname{originallemma}{Original lemma}{Original lemmas}
\makeatletter
\newlabel{lemma-gm-grothendieck-abelian}{{0FRD}{source}{}{sourcepointer-2}{}}
\newlabel{lemma-gm-grothendieck-abelian@cref}{{[originallemma][0][]0FRD}{[1][0][]source}{}{}{}}
\makeatother
\crefname{originallemma}{original lemma}{original lemmas}
\Crefname{originallemma}{Original lemma}{Original lemmas}
\makeatletter
\newlabel{remark-cone-identity}{{0FS8}{source}{}{sourcepointer-3}{}}
\newlabel{remark-cone-identity@cref}{{[originallemma][0][]0FS8}{[1][0][]source}{}{}{}}
\makeatother
\crefname{originallemma}{original lemma}{original lemmas}
\Crefname{originallemma}{Original lemma}{Original lemmas}
\makeatletter
\newlabel{section-functoriality-dg}{{0FRT}{source}{}{sourcepointer-4}{}}
\newlabel{section-functoriality-dg@cref}{{[originallemma][0][]0FRT}{[1][0][]source}{}{}{}}
\makeatother
\crefname{originallemma}{original lemma}{original lemmas}
\Crefname{originallemma}{Original lemma}{Original lemmas}
\makeatletter
\newlabel{section-functoriality-graded}{{0FR7}{source}{}{sourcepointer-5}{}}
\newlabel{section-functoriality-graded@cref}{{[originallemma][0][]0FR7}{[1][0][]source}{}{}{}}
\makeatother
\crefname{originallemma}{original lemma}{original lemmas}
\Crefname{originallemma}{Original lemma}{Original lemmas}
\makeatletter
\newlabel{section-localize-dg}{{0FRV}{source}{}{sourcepointer-6}{}}
\newlabel{section-localize-dg@cref}{{[originallemma][0][]0FRV}{[1][0][]source}{}{}{}}
\makeatother
\crefname{originallemma}{original lemma}{original lemmas}
\Crefname{originallemma}{Original lemma}{Original lemmas}
\makeatletter
\newlabel{section-localize-graded}{{0FR9}{source}{}{sourcepointer-7}{}}
\newlabel{section-localize-graded@cref}{{[originallemma][0][]0FR9}{[1][0][]source}{}{}{}}
\makeatother
\title{Stacks, Differential Graded Sheaves: lemma-dg-hull-acyclic}
\author{The Stacks Project Authors}
\date{stacks-project:sdga.tex:local-snapshot-20261012; GNU FDL 1.2 or later}
\begin{document}
\maketitle
\noindent\textit{Editorial note.} Statement span is lines 3572-3583 of sources/stacks-sdga/pinned-original-source.tex; proof span is lines 3585-3676, the complete author proof (90 body lines).\par\medskip
\noindent\textit{Editorial note.} Navigation stubs: 7 cited-but-not-included labels are supplied as non-load-bearing pointers with their pinned tags where the snapshot tag index has them.\par\medskip
\noindent\textit{Editorial note.} extra\_packages is empty on purpose: the retained preamble already loads xy with options and a second declaration clashes in the assembler wrapper.\par\medskip
\noindent\textit{Editorial note.} Text is retained verbatim, including the author's own typos and placeholders.\par\medskip
\noindent\textit{Editorial note.} Queue index 3 of queue\_sdga.json (pb=3705, stmt 3572-3583, proof 3585-3676).\par\medskip
\noindent\textit{Allowed standard prerequisites (editorial).}\begin{itemize}
\item Explicit prerequisite boundary. The cited inputs are prerequisites; the argument in the retained span is not replaced by any cited result.
\end{itemize}\medskip
\ifdefined\hypertarget\hypertarget{sourcepointer-1}{}\fi
\noindent\textit{Source reference.} Original lemma 0FSK: sources/stacks-sdga/pinned-original-source.tex, source lines 3513--3520. Cited by the retained statement or proof but not part of the collected spans; supplied only as a navigation pointer. The argument does not depend on its content.\par\medskip
\ifdefined\hypertarget\hypertarget{sourcepointer-2}{}\fi
\noindent\textit{Source reference.} Original lemma 0FRD: sources/stacks-sdga/pinned-original-source.tex, source lines 1253--1257. Cited by the retained statement or proof but not part of the collected spans; supplied only as a navigation pointer. The argument does not depend on its content.\par\medskip
\ifdefined\hypertarget\hypertarget{sourcepointer-3}{}\fi
\noindent\textit{Source reference.} Original lemma 0FS8: sources/stacks-sdga/pinned-original-source.tex, source lines 2785--2804. Cited by the retained statement or proof but not part of the collected spans; supplied only as a navigation pointer. The argument does not depend on its content.\par\medskip
\ifdefined\hypertarget\hypertarget{sourcepointer-4}{}\fi
\noindent\textit{Source reference.} Original lemma 0FRT: sources/stacks-sdga/pinned-original-source.tex, source lines 2137--2225. Cited by the retained statement or proof but not part of the collected spans; supplied only as a navigation pointer. The argument does not depend on its content.\par\medskip
\ifdefined\hypertarget\hypertarget{sourcepointer-5}{}\fi
\noindent\textit{Source reference.} Original lemma 0FR7: sources/stacks-sdga/pinned-original-source.tex, source lines 873--998. Cited by the retained statement or proof but not part of the collected spans; supplied only as a navigation pointer. The argument does not depend on its content.\par\medskip
\ifdefined\hypertarget\hypertarget{sourcepointer-6}{}\fi
\noindent\textit{Source reference.} Original lemma 0FRV: sources/stacks-sdga/pinned-original-source.tex, source lines 2248--2327. Cited by the retained statement or proof but not part of the collected spans; supplied only as a navigation pointer. The argument does not depend on its content.\par\medskip
\ifdefined\hypertarget\hypertarget{sourcepointer-7}{}\fi
\noindent\textit{Source reference.} Original lemma 0FR9: sources/stacks-sdga/pinned-original-source.tex, source lines 1014--1097. Cited by the retained statement or proof but not part of the collected spans; supplied only as a navigation pointer. The argument does not depend on its content.\par\medskip
\begin{lemma}
\label{lemma-dg-hull-acyclic}
The functors $F, G$ of Lemma \ref{lemma-dg-hull} have
the following properties. Given a graded $\mathcal{A}$-module
$\mathcal{N}$ we have
\begin{enumerate}
\item the counit $\mathcal{N} \to F(G(\mathcal{N}))$ is injective,
\item the map $\overline{\text{d}} : \mathcal{N} \to
\Coker(\mathcal{N} \to F(G(\mathcal{N})))[1]$ is an isomorphism, and
\item $G(\mathcal{N})$ is an acyclic differential graded $\mathcal{A}$-module.
\end{enumerate}
\end{lemma}

\begin{proof}
We observe that property (3) is a consequence of properties (1) and (2).
Namely, if $s$ is a nonzero local section of $F(G(\mathcal{N}))$
with $\text{d}(s) = 0$, then $s$ cannot be in the image of
$\mathcal{N} \to F(G(\mathcal{N}))$. Hence we can write the image
$\overline{s}$ of $s$ in the cokernel as
$\overline{\text{d}}(s')$ for some local section $s'$ of $\mathcal{N}$.
Then we see that $s = \text{d}(s')$ because the difference
$s - \text{d}(s')$ is still in the kernel of $\text{d}$
and is contained in the image of the counit.

\medskip\noindent
Let us write temporarily $\mathcal{A}_{gr}$, respectively $\mathcal{A}_{dg}$
the sheaf $\mathcal{A}$ viewed as a (right) graded module over itself,
respectively as a (right) differential graded module over itself.
The most important case of the lemma is to understand what
is $G(\mathcal{A}_{gr})$. Of course $G(\mathcal{A}_{gr})$ is the object
of $\textit{Mod}(\mathcal{A}, \text{d})$ representing the
functor
$$
\mathcal{M} \longmapsto
\Hom_{\textit{Mod}(\mathcal{A})}(\mathcal{A}_{gr}, F(\mathcal{M})) =
\Gamma(\mathcal{C}, \mathcal{M})
$$
By Remark \ref{remark-cone-identity} we see that this functor represented
by $C[-1]$ where $C$ is the cone on the identity of $\mathcal{A}_{dg}$.
We have a short exact sequence
$$
0 \to \mathcal{A}_{dg}[-1] \to C[-1] \to \mathcal{A}_{dg} \to 0
$$
in $\textit{Mod}(\mathcal{A}, \text{d})$ which is split by the
counit $\mathcal{A}_{gr} \to F(C[-1])$ in $\textit{Mod}(\mathcal{A})$.
Thus $G(\mathcal{A}_{gr})$ satisfies properties (1) and (2).

\medskip\noindent
Let $U$ be an object of $\mathcal{C}$. Denote
$j_U : \mathcal{C}/U \to \mathcal{C}$ the localization morphism.
Denote $\mathcal{A}_U$ the restriction of $\mathcal{A}$ to $U$.
We will use the notation $\mathcal{A}_{U, gr}$ to denote
$\mathcal{A}_U$ viewed as a graded $\mathcal{A}_U$-module.
Denote $F_U : \textit{Mod}(\mathcal{A}_U, \text{d}) \to
\textit{Mod}(\mathcal{A}_U)$ the forgetful functor and denote
$G_U$ its adjoint. Then we have the commutative diagrams
$$
\vcenter{
\xymatrix{
\textit{Mod}(\mathcal{A}, \text{d}) \ar[d]_{j_U^*} \ar[r]_F &
\textit{Mod}(\mathcal{A}) \ar[d]^{j_U^*} \\
\textit{Mod}(\mathcal{A}_U, \text{d}) \ar[r]^{F_U} &
\textit{Mod}(\mathcal{A}_U)
}
}
\quad\text{and}\quad
\vcenter{
\xymatrix{
\textit{Mod}(\mathcal{A}_U, \text{d}) \ar[r]_{F_U} \ar[d]_{j_{U!}} &
\textit{Mod}(\mathcal{A}_U) \ar[d]^{j_{U!}} \\
\textit{Mod}(\mathcal{A}, \text{d}) \ar[r]^F &
\textit{Mod}(\mathcal{A})
}
}
$$
by the construction of $j^*_U$ and $j_{U!}$ in
Sections \ref{section-functoriality-graded},
\ref{section-functoriality-dg},
\ref{section-localize-graded}, and
\ref{section-localize-dg}.
By uniqueness of adjoints we obtain $j_{U!} \circ G_U = G \circ j_{U!}$.
Since $j_{U!}$ is an exact functor, we see that the properties
(1) and (2) for the counit
$\mathcal{A}_{U, gr} \to F_U(G_U(\mathcal{A}_{U, gr}))$
which we've seen in the previous part of the proof
imply properties (1) and (2) for the counit
$j_{U!}\mathcal{A}_{U, gr} \to F(G(j_{U!}\mathcal{A}_{U, gr})) =
j_{U!}F_U(G_U(\mathcal{A}_{U, gr}))$.

\medskip\noindent
In the proof of Lemma \ref{lemma-gm-grothendieck-abelian}
we have seen that any object
of $\textit{Mod}(\mathcal{A})$ is a quotient of a direct sum
of copies of $j_{U!}\mathcal{A}_{U, gr}$. Since $G$ is a left
adjoint, we see that $G$ commutes with direct sums. Thus
properties (1) and (2) hold for direct sums of objects
for which they hold. Thus we see that every object $\mathcal{N}$
of $\textit{Mod}(\mathcal{A})$ fits into an exact sequence
$$
\mathcal{N}_1 \to \mathcal{N}_0 \to \mathcal{N} \to 0
$$
such that (1) and (2) hold for $\mathcal{N}_1$ and $\mathcal{N}_0$.
We leave it to the reader to deduce (1) and (2) for
$\mathcal{N}$ using that $G$ is right exact.
\end{proof}

\end{document}
